\section{Обработка пропусков в дереве}

Как мы уже знаем, $NaN$ (пропуски) в данных это сама по себе довольно насущная проблема. Однако у решающих деревьев есть один довольно необычный метод борьбы с этим "--- \textit{метод суррогатных предикатов.}

Идея метода суррогатных предикатов, первоначально предложенного Брейманом и его коллегами в алгоритме CART, заключается в том, что для каждого основного правила разбиения в узле дерева мы находим дополнительные, резервные правила, которые используют другие признаки, но дают максимально похожее разбиение данных.

Представьте себе врача, который ставит диагноз пациенту. Если невозможно измерить основной показатель (например, уровень определенного гормона в крови), опытный врач будет опираться на дополнительные симптомы и анализы, которые обычно коррелируют с этим основным показателем. Аналогичным образом, суррогатные предикаты выступают в роли таких «дополнительных симптомов», позволяя принять обоснованное решение даже при отсутствии части информации.

Формально, для каждого основного правила разбиения вида $f(x) \leq c$ мы находим набор суррогатных правил $g_i(x) \leq d_i$, которые максимально близко воспроизводят разбиение, создаваемое основным правилом. Эти суррогатные правила упорядочиваются по степени их соответствия основному разбиению, образуя иерархическую систему резервных вариантов.

Обычно, если не удаётся установить суррогатные предикаты по тем или иным причинам, то появляется волшебная фраза \textit{вносим свои эвристики}. На самом деле это настолько редкая ситуация, что до такого крайне редко доходит (особенно если данные всё"=таки обрабатываются, и поступают в модель без пропусков).

\subsection{Оценка качества суррогатных правил}

Ключевым аспектом метода суррогатных предикатов является количественная оценка того, насколько хорошо каждое суррогатное разбиение приближает основное. Для этого используется специальная мера соответствия, которая вычисляет долю объектов, попадающих в одинаковые подмножества при применении как основного, так и суррогатного правила.

Если обозначить через $S$ разбиение, создаваемое основным правилом, а через $S'$ — разбиение суррогатного предиката, то мера соответствия вычисляется как отношение числа объектов, которые попали в одинаковые группы при обоих разбиениях, к общему числу объектов. Математически это выражается формулой:
\[
\text{совпадение}(S, S') = \frac{\text{число объектов с одинаковым назначением в } S \text{ и } S'}{\text{общее число объектов}}
\]

На практике суррогатные предикаты сортируются по убыванию этой меры соответствия, формируя упорядоченный список запасных правил. Когда встречается объект с пропуском в основном признаке, алгоритм последовательно перебирает суррогатные предикаты, начиная с наиболее соответствующего, до тех пор, пока не найдет правило, для которого у данного объекта есть значение.
